\subsection{根系的直和}

设 \(V\) 是 \(k\) 上的向量空间，且是向量空间族 \((V_i)_{1 \leqslant i \leqslant r}\) 的直和。将 \(V^*\) 与 \(V_i^*\) 的直和等同。对所有 \(i\)，设 \(R_i\) 是 \(V_i\) 中的根系。则 \(R = \bigcup_i R_i\) 是 \(V\) 中的根系，其逆根系为 \(R^- = \bigcup_i R_i^-\)；从 \(R\) 到 \(R^-\) 的典范双射对每个 \(i\) 都扩张了从 \(R_i\) 到 \(R_i^-\) 的典范双射。集合 \(R\) 称为根系 \(R_i\) 的直和。设 \(\alpha \in R_i\)。若 \(j \neq i\)，则 \(\alpha^-\) 的核包含 \(V_j\)，所以 \(s_\alpha\) 在 \(V_j\) 上诱导恒等映射；另一方面，\(k\alpha \subseteq V_i\)，所以 \(s_\alpha\) 保持 \(V_i\) 不变。这些说明表明，\(W(R)\) 可与 \(\prod_{i=1}^{r} W(R_i)\) 等同。

称根系 R 是不可约的，当且仅当 \(R \neq \varnothing\)，且 R 不是两个非空根系的直和。

命题 5. 设 V 是 k 上的向量空间，且 V 是向量空间 \(V_{1}, \ldots, V_{r}\) 的直和。设 R 是 V 中的根系。置 \(R_{i} = R \cap V_{i}\)。以下三个条件等价：

\begin{enumerate}
\def\labelenumi{(\roman{enumi})}
\item
  \(V_{i}\) 在 W(R) 下不变；
\item
  \(R \subseteq V_{1} \cup V_{2} \cup \cdots \cup V_{r};\)
\item
  对所有 \(i\)，\(\mathbb{R}_i\) 是 \(\mathrm{V}_i\) 中的根系，且 \(\mathbb{R}\) 是 \(\mathbb{R}_i\) 的直和。
\item
  \(\Longrightarrow\) (i)：这由本号开头所述得到。
\item
  \(\Longrightarrow\) (ii)：假设 \(V_{i}\) 在 \(W(R)\) 下不变。设 \(\alpha \in R\)，并设 H 是 \(\alpha\) 的核。由 Chap. V, § 2, no. 2 的 Prop. 3，每个 \(V_{i}\) 都是 H 的一个子空间与 \(k\alpha\) 的一个子空间之和。因此某个 \(V_{i}\) 包含 \(k\alpha\)，所以 \(\alpha \in V_{1} \cup V_{2} \cup \cdots \cup V_{r}\)。
\item
  \(\Longrightarrow\) (iii)：若条件 (ii) 满足，则 \(\mathbf{R}_i\) 生成 \(\mathbf{V}_i\) 对所有 \(i\) 成立，所以 \(\mathbf{R}_i\) 是 \(\mathbf{V}_i\) 中的根系（命题 4）。显然 \(\mathbf{R}\) 是 \(\mathbf{R}_i\) 的直和。
\end{enumerate}

推论。设 R 是 V 中的根系。以下条件等价：

\begin{enumerate}
\def\labelenumi{(\roman{enumi})}
\item
  R 不可约：
\item
  W(R)-模 V 是单的；
\item
  W(R)-模 V 是绝对单的。
\item
  \(\Longleftrightarrow\) (i)：这由命题 5 和 Maschke 定理（Chap. V, Appendix, Prop. 2）得到。
\item
  \(\Longleftrightarrow\) (ii)：这由 Chap. V, § 2, no. 1 的 Prop. 1 得到。
\end{enumerate}

命题 6. V 中的每个根系 R 都是一个不可约根系族 \((\mathrm{R}_{i})_{i\in I}\) 的直和，并且除指标集的一个双射外是唯一的。

\(R_{i}\) 的存在性由对 Card R 归纳证明：若 R 非空且可约，则 R 是两个根系 \(R'\)、\(R''\) 的直和，满足 \(Card R' < Card R\)、\(Card R'' < Card R\)，并且归纳假设适用于 \(R'\) 和 \(R''\)。为证明唯一性，只需证明：若 R 是 \(R'\) 与 \(R''\) 的直和，则每个 \(R_i\) 必然包含于 \(R'\) 或 \(R''\) 中。设 \(V', V'', V'_i, V''_i\) 是 V 中分别由 \(R', R'', R' \cap R_i, R'' \cap R_i\) 生成的向量子空间。由于和 \(V' + V''\) 是直和，和 \(V'_i + V''_i\) 也是直和。由于 \(R_i \subseteq R' \cup R''\)，\(R_i\) 是根系 \(R' \cap R_i\) 与 \(R'' \cap R_i\) 的直和；因此 \(R' \cap R_i = \varnothing\) 或 \(R'' \cap R_i = \varnothing\)，这就证明了所述断言。

\(R_{i}\) 称为 R 的不可约分支。对于任意非零标量 \(\lambda_{i}\)，\(\lambda_{i}R_{i}\) 的并集是 V 中的根系，其逆根系是 \(\lambda_{i}^{-1}R_{i}\) 的并集，并且其 Weyl 群为 \(\mathrm{W}(\mathrm{R})\)。

命题 7. 设 R 是 V 中的根系，(R \(_{i}\) ) 是其不可约分支族，V \(_{i}\) 是由 R \(_{i}\) 生成的 V 的向量子空间，B 是命题 3 中定义的 V 上的不变对称双线性型，B' 是 V 上在 W(R) 下不变的对称双线性型。则各 V \(_{i}\) 关于 B' 两两正交，并且对所有 i，B 和 B' 在 V \(_{i}\) 上的限制成比例。

若 \(v_{i} \in \mathrm{V}_{i}\)、\(v_{j} \in \mathrm{V}_{j}\)、\(i \neq j\)，且 \(w \in \mathrm{W}(\mathrm{R}_j)\)，则

\[
\mathrm{B} ^ {\prime} (v _ {i}, w (v _ {j})) = \mathrm{B} ^ {\prime} (v _ {i}, v _ {j}),
\]

这表明 \(w(v_{j}) - v_{j}\) 与 \(v_{i}\) 关于 \(B'\) 正交。由于 \(V_{j}\) 对 \(\mathrm{W}(\mathrm{R}_{j})\) 不可约，它由 \(w(v_{j}) - v_{j}\) 生成，因此它与 \(V_{i}\) 正交。

B 和 B' 在各个 \(V_{i}\) 上的限制成比例这一事实由 Chap. V, § 2, no. 1 的 Prop. 1 得到。

注。取一个 \(V_{R}\) 上在 \(W(R)\) 下不变的内积。于是可以谈论根的长度以及两个根之间的角。命题 7 表明，只要两个根属于 R 的同一不可约分支，这个角与内积的选择无关，两个根的长度之比也与内积的选择无关。

